\section{Singer convolution and the negation budget}
This section computes all projective-plane line disjunctions by cyclic
convolution and accounts for the negations needed to share that computation.
Let $N=q^2+q+1$, and let $C_s(q)$ be the least size of a fan-in-two
AND/OR/NOT circuit computing all line ORs of $\operatorname{PG}(2,q)$
with at most $s$ NOT gates. Outputs and fan-out are free.

\subsection{Singer coordinates}
The cyclic incidence description has the following direct finite-field
proof; historically it is Singer's construction \cite{Singer1938}.
Put $K=\mathbb F_{q^3}$ and identify projective points with
$K^\times/\mathbb F_q^\times$, a cyclic group of order $N$.
Choose a generator $\alpha$ of $K^\times$.
The lines are the projectivized kernels of
$x\mapsto\operatorname{Tr}_{K/\mathbb F_q}(\beta x)$, for
$\beta\in K^\times$. Indeed, trace is a nonzero linear functional:
the nonzero polynomial $T+T^q+T^{q^2}$ has degree less than $q^3$,
so it cannot vanish everywhere on $K$. Multiplication by nonzero
$\beta$ is invertible; hence the trace pairing is nondegenerate and
exhausts the three-dimensional dual space.
Two nonzero functionals have the same kernel exactly when their
ratio is a scalar in $\mathbb F_q^\times$.

Let $D=\{d\in\mathbb Z/N:\operatorname{Tr}(\alpha^d)=0\}$.
The trace kernel has $q^2-1$ nonzero vectors, so $|D|=q+1$.
Multiplication by $\alpha$ acts regularly on points and lines,
and the line sets are precisely $D+w$, $w\in\mathbb Z/N$.
Distinct lines meet once; in particular
$|D\cap(D+w)|=1$ for $w\ne0$.

\subsection{Convolution and a shared input inverter}

In Singer coordinates the outputs are
$y_w=\bigvee_{d\in D}x_{w+d}$ on $\mathbb Z/N$, with $|D|=q+1$.
For $L=\lceil\log_2(q+2)\rceil$ and $R=2^L$, form
\[
 X=\sum_{i=0}^{N-1}x_iR^i,\qquad B=\sum_{d\in D}R^{N-d}.
\]
Every base-$R$ coefficient of $XB$ is at most $q+1<R$, and the
coefficients at positions $w$ and $w+N$ sum to the number of true
inputs on line $w$. One integer multiplication followed by short
additions and nonzero tests gives
\[
 \mathcal S_N=O(N\log^2N\log\log N)
\]
gates. The network form of the multiplication bound is explicitly
part of Sch\"onhage--Strassen \cite[Summary]{SchonhageStrassen1971}.
The factors have at most $(N+1)L$ bits; their logical-network
multiplication bound is $O(t\log t\log\log t)$ for $t$-bit integers.
The $N$ block additions and nonzero tests cost $O(NL)$ further
gates. An arbitrary binary Boolean gate can be replaced by a
constant number of gates in the AND/OR/NOT basis.

Beals--Nishino--Tanaka \cite[Theorem~1.1, p.~1335;
Corollary~1.2, p.~1336]{BealsNishinoTanaka1998}
give an $m$-input inverter of size $O(m\log m)$ with
$b(m)=\lceil\log_2(m+1)\rceil$ negations, and the simulation
$C^{b(m)}(F)\leq2C(F)+O(m\log m)$ for arbitrary systems $F$.
The dual-rail implementation is explicit. Supply each input
$x_i$ and $\neg x_i$ using one shared $m$-input inverter.
If $a^+,a^-$ and $b^+,b^-$ are complementary pairs, implement an AND
by $(a^+\wedge b^+,a^-\vee b^-)$ and an OR by
$(a^+\vee b^+,a^-\wedge b^-)$.
A NOT merely interchanges two references.
Every original AND or OR costs at most two gates, original NOT
gates cost none, and positive wires at all outputs compute the
whole system. Thus the negation count is solely that of the
shared input inverter, independently of the number of outputs.
Explicitly counted constant gates contribute at most a fixed
additive allowance.

Fischer's older simulation has the same structure but a larger
inverter bound. His Theorem~16, printed p.~8 of the MIT memorandum
\cite{Fischer1975}, gives
\[
 NC_{b(m)}(F)\le 2C^\Delta(F)+O(m^2\log^2m)
              \le 6C(F)+O(m^2\log^2m),
\]
where $C^\Delta$ uses AND/OR/NOT and the second $C$ uses the full
unary/binary basis. The sharper additive term used here belongs
to Beals--Nishino--Tanaka.

\subsection{An endpoint refinement}


\begin{proposition}
For every prime power $q\geq2$,
\[
 C_{\lceil\log_2(q^2)\rceil}(q)
 \leq O(N\log^2N\log\log N).
\]
In particular, if $q=2^r$, the same bound holds with exactly the
budget $s=2r=\lfloor\log_2N\rfloor$.
\end{proposition}
\begin{proof}
Leave unprocessed a set $R_0$ consisting of one line and one point
outside it. Its complement has $m=q^2-1$ points. Use the fast
convolution circuit with all inputs in $R_0$ fixed to zero, and
apply the inverter simulation only to the $m$ remaining inputs.
This costs $2\mathcal S_N+O(m\log m)$ gates and
$b(m)=\lceil\log_2(q^2)\rceil$ negations.

The contributions from $R_0$ require $q$ OR gates on its complete
line, and one OR gate on each of the $q+1$ lines through the extra
point. All other contributions are single input wires.
Joining these with the fast contributions uses $N-1$ further
OR gates. The added cost is therefore $N+2q$, and no new negation
is used.
\end{proof}
More generally, choosing any $1\leq m\leq\min(N,2^s-1)$ inputs for the fast
part gives
\[
 C_s(q)\leq
 \min\{Nq,\ 2\mathcal S_N+O(m\log(m+1))
 +(N-m)(q+1)-\lambda(R_0)\},
\]
where $\lambda(R_0)$ counts the complete lines in the complementary
point set. A constant additive allowance handles explicit constant
gates. This is an upper envelope, not a determination of $C_s$.

If a fixed $c>0$ satisfied $C_s(q)\geq c^s qN$ for all prime-power
$q$ and all integer $s\leq\log_2N$, the proposition applied at
$q=2^r$, $s=2r$ would imply
$(2c^2)^r=O(r^2\log r)$. Hence $c\leq2^{-1/2}$.
This argument respects the integer budget; the generic simulation
at $\lceil\log_2(N+1)\rceil$ alone does not justify comparison with a
lower bound assumed only through $\lfloor\log_2N\rfloor$.
